\chapter{التكامل على قطعة}\label{ch:b1:integration}

قام تكامل مجلّد الثانوي على مساحات مأخوذة
بالحدس. ويبني هذا الفصل \omterm{thm:b1:integration:def}{التكامل}: أولًا من أجل الدوال الدرجية،
حيث يكون \omterm{thm:b1:integration:def}{التكامل} مجموعًا منتهيًا، ثم من أجل الدوال \omterm{def:b1:continuity:continuous}{المتصلة} (\omterm{def:b1:continuity:continuous}{والمتصلة}
قطعةً قطعة) بالتقريب المنتظم --- وهو الموضع الذي تكسب فيه
مبرهنة هاينه (\cref{thm:b1:continuity:heine}) قوتَها.
ثم تربط المبرهنة الأساسية في \omterm{thm:b1:integration:def}{التكامل} البناءَ
بالدوال الأصلية، وتربطه \omterm{thm:b1:integration:riemann}{مجاميع ريمان} بالمتوسطات المتقطّعة.

وفي كل ما يأتي، $a < b$ عددان حقيقيان.

\section{الدوال الدرجية}

\begin{definition}\label{def:b1:integration:step}
تكون $\varphi \colon \intcc{a}{b} \to \R$ \emph{دالة
درجية}\index{دالة درجية} إذا وُجدت تجزئة $a = x_0
< x_1 < \dots < x_n = b$ تكون $\varphi$ عليها ثابتة، مساوية
$c_i$، على كل \omterm{prop:b1:reals:intervals}{فترة} \omterm{def:b1:topology:open}{مفتوحة} $\intoo{x_{i-1}}{x_i}$ (والقيم عند
العقد غير مقيَّدة). وتكاملها هو
\[
\int_a^b \varphi = \sum_{i=1}^{n} c_i\,(x_i - x_{i-1}),
\]
وهو مستقل عن التجزئة المختارة (لطّف تجزئتين
بتجزئتهما المشتركة: فلا يتغير أيّ طرف بالتلطيف).
\end{definition}

\begin{proposition}\label{prop:b1:integration:stepprops}
على الدوال الدرجية، يكون \omterm{thm:b1:integration:def}{التكامل} خطيًا ومتزايدًا ($\varphi \leq
\psi \implies \int\varphi \leq \int\psi$)، ويحقق علاقة
شال $\int_a^b = \int_a^c + \int_c^b$ من أجل $a < c < b$.
\end{proposition}

\begin{proof}
المحرّك هو \emph{الحفظ بالتلطيف}، المذكور في
التعريف: فإدخال عقدة إضافية واحدة $t \in \intoo{x_{i-1}}{x_i}$
في تجزئة يستبدل بالحد $c_i(x_i - x_{i-1})$ المقدارَ
$c_i(t - x_{i-1}) + c_i(x_i - t)$ --- وهو العدد نفسه --- ومنه فلا
يتغير \omterm{thm:b1:integration:def}{التكامل} بأيّ تلطيف منته. والآن خذ
$\varphi$ بالتجزئة $\sigma$ و $\psi$ بالتجزئة
$\sigma'$: فعلى التلطيف المشترك $\sigma \cup \sigma'$ تكون كلتاهما
دالتين درجيتين بالعقد \emph{نفسها}، وعلى كل قطعة تكون
$\varphi + \lambda\psi$ ثابتة تساوي $c_i + \lambda
d_i$: ومنه تُردّ الخطية إلى خطية المجاميع المنتهية. والتزايد:
يعطي $c_i \leq d_i$ على كل قطعة أن $\sum c_i \Delta_i \leq \sum
d_i \Delta_i$ (لأن الأطوال $\Delta_i \geq 0$). وشال: أدخل
العقدة $c$ واقسم المجموع عندها.
\end{proof}

\section{تكامل دالة متصلة}

\begin{theorem}[التقريب المنتظم]\label{thm:b1:integration:approx}
لتكن $f$ \omterm{def:b1:continuity:continuous}{متصلة} على $\intcc{a}{b}$. من أجل كل $\varepsilon > 0$
توجد دالتان درجيتان $\varphi, \psi$ تحققان
\[
\varphi \leq f \leq \psi
\qquad\text{و}\qquad
\psi - \varphi \leq \varepsilon \text{ على } \intcc{a}{b}.
\]
\end{theorem}

\begin{proof}
حسب مبرهنة هاينه (\cref{thm:b1:continuity:heine})، تكون $f$ \omterm{def:b1:continuity:uniform}{متصلة
انتظامًا}: اختر $\delta$ من أجل $\varepsilon$، وتجزئةً
خطوتها $< \delta$ (متساوية التباعد مثلًا، بعدد $n > \frac{b -
a}{\delta}$ قطعة). وعلى كل قطعة \omterm{def:b1:topology:closed}{مغلقة} $\intcc{x_{i-1}}{x_i}$، تبلغ $f$
قيمة صغرى $m_i$ وقيمة عظمى $M_i$
(\cref{thm:b1:continuity:evt})، و $M_i - m_i \leq \varepsilon$
(لأن النقطتين الحدّيتين على بعد $\delta$). عرّف $\varphi = m_i$
و $\psi = M_i$ على $\intoo{x_{i-1}}{x_i}$ (و $\varphi = \psi = f$
عند العقد).
\end{proof}

\begin{theorem}[تعريف التكامل]\label{thm:b1:integration:def}
لتكن $f$ \omterm{def:b1:continuity:continuous}{متصلة} على $\intcc{a}{b}$. العددان
\[
I_-(f) = \sup\Bigl\{\int_a^b \varphi : \varphi \text{ درجية},\
\varphi \leq f\Bigr\},
\qquad
I_+(f) = \inf\Bigl\{\int_a^b \psi : \psi \text{ درجية},\ \psi \geq
f\Bigr\}
\]
متساويان؛ وقيمتهما المشتركة هي
\emph{التكامل}\index{تكامل} $\int_a^b f$ (ويُكتب كذلك $\int_a^b
f(t)\,\dd t$). وهو يوافق المفهوم السابق على الدوال
الدرجية، ويمتدّ إلى الدوال \omterm{def:b1:continuity:continuous}{المتصلة} قطعةً قطعة
بقطع $\intcc{a}{b}$ عند الانقطاعات (وشال تعريفًا
هناك).
\end{theorem}

\begin{proof}
المجموعتان غير خاليتين (لأن $f$ محدودة) وكل تكامل درجي أدنى
هو $\leq$ كل تكامل أعلى (بالرتابة على الدوال الدرجية): ومنه
$I_-(f) \leq I_+(f)$. وحسب \cref{thm:b1:integration:approx}، يوجد من أجل كل
$\varepsilon$ زوجٌ يحقق $\int\psi - \int\varphi \leq
\varepsilon(b - a)$: ومنه يُحصر الحدّان الأعلى والأدنى معًا، $I_- =
I_+$.
\end{proof}

\begin{example}[متصلة قطعةً قطعة، بلا دراما]\label{ex:b1:integration:piecewise}
\omterm{thm:b1:reals:floor}{دالة الجزء الصحيح} على $\intcc{0}{3}$ \omterm{def:b1:integration:step}{دالة درجية}
متنكّرة: فبالقطع عند قفزاتها،
\[
\int_0^3 \lfloor t \rfloor\,\dd t
= \int_0^1 0 + \int_1^2 1 + \int_2^3 2 = 0 + 1 + 2 = 3 ,
\]
والقيم \emph{عند} نقطتَي القفز $1, 2$ لا شأن لها:
فتغيير دالة عند عدد منته من النقاط لا يغيّر أيّ تكامل
(لأن الدوال الدرجية المؤطِّرة لا تتأثر). وهذا هو كل
مضمون التمديد إلى «\omterm{def:b1:continuity:continuous}{المتصلة} قطعةً قطعة»: اقطع عند
الانقطاعات المنتهية العدد، وكامل كل قطعة \omterm{def:b1:continuity:continuous}{متصلة}، ثم
اجمع --- أي شال تعريفًا.
\end{example}

\begin{example}[التعريف يحسب، مرة واحدة]\label{ex:b1:integration:defcompute}
لتكن $f(x) = x$ على $\intcc{0}{1}$ واقطع إلى $n$ قطعة متساوية.
وأفضل الدوال الدرجية الثابتة على القطع هي $\varphi =
\frac{k-1}{n}$ و $\psi = \frac kn$ على القطعة $k$، مع
\[
\int_0^1 \varphi = \sum_{k=1}^{n} \frac{k-1}{n}\cdot\frac1n
= \frac{n-1}{2n},
\qquad
\int_0^1 \psi = \sum_{k=1}^{n} \frac{k}{n}\cdot\frac1n
= \frac{n+1}{2n} .
\]
وكل تكامل أدنى هو $\leq I_-(f) \leq I_+(f) \leq$ كل
تكامل أعلى، ومنه $\frac{n-1}{2n} \leq I_-(f) \leq I_+(f) \leq
\frac{n+1}{2n}$ من أجل كل $n$: فيُحصر كلاهما على $\frac12$، و
$\int_0^1 x\,\dd x = \frac12$ مباشرةً من التعريف. والفكرة
النافذة: هذه أول وآخر مرة نكامل فيها
من التعريف --- إذ تستبدل المبرهنة الأساسية أدناه
بكل هذه الحسابات بحثًا واحدًا عن دالة أصلية، وهو
كل الفائدة الاقتصادية لهذا الفصل.
\end{example}

\begin{theorem}[الخصائص]\label{thm:b1:integration:props}
من أجل $f, g$ متصلتين (أو متصلتين قطعةً قطعة) على $\intcc{a}{b}$ و
$\lambda \in \R$:
\begin{enumerate}
  \item الخطية: $\int (f + \lambda g) = \int f + \lambda \int g$؛
  \item الرتابة: $f \leq g \implies \int_a^b f \leq \int_a^b g$؛
        و $\bigl|\int_a^b f\bigr| \leq \int_a^b \abs f \leq (b -
        a)\, \sup\abs f$؛
  \item شال: $\int_a^b = \int_a^c + \int_c^b$ (مع
        الاصطلاح $\int_b^a = -\int_a^b$، وهو صحيح من أجل أيّ ترتيب
        للحدّين)؛
  \item الإيجابية القطعية: إذا كانت $f$ \omterm{def:b1:continuity:continuous}{متصلة} و $f \geq 0$ و
        $\int_a^b f = 0$، فإن $f = 0$ في كل مكان على
        $\intcc{a}{b}$.
\end{enumerate}
\end{theorem}

\begin{proof}
تمرّ البنود (1)--(3) من الدوال الدرجية إلى النهاية عبر
التعريف \omterm{def:b1:reals:bounds}{بالحد الأعلى} والأدنى. وتستحق الخطية التفاصيل مرة واحدة: فمن أجل
$\varepsilon > 0$ معطى، أطّر $\varphi_f \leq f \leq \psi_f$ و
$\varphi_g \leq g \leq \psi_g$ بفجوتين $\leq \varepsilon$
(\cref{thm:b1:integration:approx}). ومن أجل $\lambda \geq 0$، يكون
$\varphi_f + \lambda\varphi_g \leq f + \lambda g \leq \psi_f +
\lambda\psi_g$ تأطيرًا بدوال درجية بفجوة $\leq
(1 + \lambda)\varepsilon$، وتساوي تكاملاته الدرجية $\int
\varphi_f + \lambda\int\varphi_g$ وهكذا
(\cref{prop:b1:integration:stepprops}): ومنه يحصر جعلُ $\varepsilon \to
0$ المقدارَ $\int(f + \lambda g)$ على $\int f + \lambda\int g$.
وأمّا من أجل $\lambda < 0$، فالضرب في $\lambda$ \emph{يقلب}
تأطير $g$ --- إذ تصير \omterm{def:b1:integration:step}{الدالة الدرجية} الدنيا للمقدار $\lambda g$ هي
$\lambda \psi_g$ --- ويجري الحصر نفسه بتبادل
الدورين. والحاصر $\abs{\int f} \leq \int \abs f$ يأتي من
$-\abs f \leq f \leq \abs f$ ومن الرتابة.

(4) بالنقيض: إذا كان $f(x_0) = m > 0$، فإن \omterm{def:b1:continuity:continuous}{الاتصال} يعطي
\omterm{prop:b1:reals:intervals}{فترة} جزئية طولها $\eta > 0$ تحقق عليها $f \geq \frac m2$؛ و
\omterm{def:b1:integration:step}{الدالة الدرجية} التي تساوي $\frac m2$ هناك و $0$ فيما عداه هي $\leq f$،
ومنه $\int f \geq \frac{m\eta}{2} > 0$.
\end{proof}

\begin{example}[شال في العمل: تكاملات ذات قيم مطلقة]\label{ex:b1:integration:chasleswork}
لمكاملة قيمة مطلقة، اقطع حيث تتغير الإشارة.
\[
\int_0^2 \abs{x - 1}\,\dd x
= \int_0^1 (1 - x)\,\dd x + \int_1^2 (x - 1)\,\dd x
= \frac12 + \frac12 = 1 ,
\]
و، بقطع $\intcc{0}{2\pi}$ عند $\pi$:
\[
\int_0^{2\pi} \abs{\sin t}\,\dd t
= \int_0^{\pi} \sin t\,\dd t - \int_{\pi}^{2\pi} \sin t\,\dd t
= 2 + 2 = 4 ,
\]
بينما $\int_0^{2\pi} \sin t\,\dd t = 0$: فالتلاشي حقيقي،
ولهذا تحمل \omterm{def:b1:logic:statement}{عبارة} الإيجابية القطعية
(\cref{thm:b1:integration:props}~(4)) الفرضَ $f
\geq 0$ --- فبدونه، لا يبرهن \omterm{thm:b1:integration:def}{التكامل} المنعدم على شيء
عن $f$. والفكرة النافذة: يقيس $\int \abs f$
\emph{المساحة}، ويقيس $\int f$ \emph{الميزان الموقَّع}؛ و
المتراجحة $\abs{\int f} \leq \int\abs f$ سجلٌّ مضبوط
لما يستطيع التلاشي تدميره.
\end{example}

\section{المبرهنة الأساسية في التكامل}

\begin{theorem}[المبرهنة الأساسية في التكامل]\label{thm:b1:integration:ftc}
لتكن $f$ \omterm{def:b1:continuity:continuous}{متصلة} على \omterm{prop:b1:reals:intervals}{فترة} $I$ وليكن $a \in I$. الدالة
\[
F(x) = \int_a^x f(t)\, \dd t
\]
من الصنف $C^1$ على $I$، مع $F' = f$: ومنه فلكل دالة \omterm{def:b1:continuity:continuous}{متصلة}
على \omterm{prop:b1:reals:intervals}{فترة} دوالُّ أصلية. ونتيجةً لذلك، من أجل \emph{أيّ}
دالة أصلية $G$ للمقدار $f$:\index{المبرهنة الأساسية في التكامل}
\[
\int_a^b f(t)\,\dd t = G(b) - G(a) .
\]
\end{theorem}

\begin{proof}
ثبّت $x_0 \in I$ و $\varepsilon > 0$؛ يعطي \omterm{def:b1:continuity:continuous}{الاتصال} عند $x_0$
عددًا $\delta$ يحقق $\abs{f(t) - f(x_0)} \leq \varepsilon$ من أجل $\abs{t -
x_0} \leq \delta$. ومن أجل $0 < \abs{h} \leq \delta$ (و $x_0 + h \in
I$)، تعطي شال
\[
\frac{F(x_0 + h) - F(x_0)}{h} - f(x_0)
= \frac 1h \int_{x_0}^{x_0+h} \bigl(f(t) - f(x_0)\bigr)\dd t ,
\]
وقيمته المطلقة هي $\leq \frac{1}{\abs h}\cdot \abs h\,
\varepsilon = \varepsilon$ (بالحاصر (2)، الصحيح من أجل أيّ ترتيب
للحدّين). ومنه $F'(x_0) = f(x_0)$؛ و $F' = f$ \omterm{def:b1:continuity:continuous}{متصلة}: ومنه فالدالة $F$
من الصنف $C^1$. وإذا كان $G' = f$ كذلك، فإن $(G - F)' = 0$ على \omterm{prop:b1:reals:intervals}{الفترة}، ومنه $G =
F + c$ (\cref{cor:b1:derivative:monotone})، و $G(b) - G(a) = F(b)
- F(a) = \int_a^b f$.
\end{proof}

\begin{example}[التناظر قبل الحساب]\label{ex:b1:integration:parity}
على \omterm{prop:b1:reals:intervals}{فترة} متناظرة، تقوم الزوجية بالعمل: فإذا كانت $f$ فردية،
يرسل التعويض $t \mapsto -t$ المقدارَ $\int_{-a}^{0} f$ إلى
$-\int_0^a f$، ومنه
\[
\int_{-a}^{a} f(t)\,\dd t = 0 ;
\qquad\text{وإذا كانت $f$ زوجية،}\quad
\int_{-a}^{a} f = 2\int_0^a f .
\]
ومنه $\int_{-1}^{1} \frac{t^3\cos t}{1 + t^4}\,\dd t = 0$ دون
أيّ دالة أصلية تلوح (لأن المكامَل فردي)، و
$\int_{-\pi}^{\pi} t^2\cos t\,\dd t = 2\int_0^\pi t^2\cos t\,
\dd t$. فتحقق من التناظر قبل أن تمدّ يدك إلى التقنيات: فأسرع
تكامل هو الذي لا يُحسب قط.
\end{example}

\begin{example}[التعرّف على مشتقة بالنظر]\label{ex:b1:integration:halfangle}
احسب $\displaystyle\int_0^{\pi/2} \frac{\dd x}{1 + \cos x}$.
تحوّل متطابقة نصف الزاوية $1 + \cos x = 2\cos^2\frac x2$ المكامَلَ
إلى $\frac{1}{2}\bigl(1 + \tan^2\frac
x2\bigr)$، وهو بالضبط \omterm{def:b1:derivative:def}{مشتقة} $\tan\frac x2$:
\[
\int_0^{\pi/2} \frac{\dd x}{1 + \cos x}
= \Bigl[\tan\frac x2\Bigr]_0^{\pi/2} = \tan\frac\pi4 = 1 .
\]
ولم يُحتَج إلى أيّ آلة تعويض --- بل إلى منعكس
قراءة المكامَل مشتقةً لأحدهم فقط، والمبرهنة
الأساسية تفعل الباقي. (وأمّا الأداة المنهجية وراء مثل هذه
التكاملات المثلثية، أي التعويض $t = \tan\frac x2$،
فتنتمي إلى العدّة المعيارية المبنية من
\cref{thm:b1:integration:parts}~(2).)
\end{example}

\begin{example}[دوال معرَّفة بتكاملات]\label{ex:b1:integration:definedby}
تصنع المبرهنة الأساسية دوالًّ. لتكن
\[
F(x) = \int_0^x \eu^{-t^2}\,\dd t .
\]
لا تركيبة من الدوال الكلاسيكية مشتقتُها
$\eu^{-t^2}$ (وهي مبرهنة لليوفيل، مقبولة هنا)؛ ومع ذلك توجد $F$،
وهي من الصنف $C^1$ مع $F'(x) = \eu^{-x^2} > 0$، ومتزايدة قطعًا، وفردية
(بالتعويض $t \mapsto -t$)، ومحدودة: فمن أجل $x \geq 1$،
\[
F(x) - F(1) = \int_1^x \eu^{-t^2}\dd t
\leq \int_1^x \eu^{-t}\dd t \leq \eu^{-1} ,
\]
ومنه $F \leq F(1) + \eu^{-1} \leq 1 + \eu^{-1}$. (والنهاية المضبوطة،
$\frac{\sqrt\pi}{2}$، تُحسب بالتكاملات المضاعفة في
مجلّد السنة الثالثة.) وقاعدة السلسلة من أجل حدود متحركة: $\frac{\dd}{\dd
x}\int_x^{x^2} \eu^{-t^2}\dd t = 2x\,\eu^{-x^4} - \eu^{-x^2}$.
والفكرة النافذة: \omterm{thm:b1:integration:def}{التكامل} \emph{ينشئ} دوالًّ جديدة
من قديمة، بكل خصائصها مقروءةً من
المكامَل --- فالدالة الأصلية التي لا تستطيع كتابتها تبقى
دالة تسيطر عليها سيطرة تامة.
\end{example}

\begin{example}[التقدير دون التقويم]\label{ex:b1:integration:estimate}
التكاملات $R_n = \int_0^1 \frac{t^n}{1 + t}\,\dd t$ ليس لها أيّ صورة
مغلقة لطيفة، ومع ذلك تثبّتها الرتابة بدقة: فعلى
$\intcc{0}{1}$، $\frac12 \leq \frac{1}{1+t} \leq 1$، ومنه
\[
\frac{1}{2(n+1)} = \frac12\int_0^1 t^n\,\dd t
\;\leq\; R_n \;\leq\; \int_0^1 t^n \,\dd t = \frac{1}{n+1} :
\]
أي رتبة التلاشي المضبوطة ($R_n \sim$ مضاعف للمقدار $\frac1n$،
وفي الواقع $R_n \sim \frac{1}{2n}$) بسطرين ودون أيّ
دالة أصلية. وتجري مسألتا نهاية الأسبوع في هذا الفصل والذي
يليه على مثل هذه التأطيرات بالضبط --- فأول منعكس عند المحلل
أمام تكامل ينبغي أن يكون \emph{حُدَّه}، وعندئذ
فقط، إن لزم، احسبه.
\end{example}

\begin{example}[القيم المتوسطة]\label{ex:b1:integration:average}
\emph{متوسط} دالة \omterm{def:b1:continuity:continuous}{متصلة} $f$ على $\intcc{a}{b}$ هو
$\frac{1}{b-a}\int_a^b f$. ومن أجل \omterm{def:b1:functions:arc}{قوس الجيب}:
\[
\frac{1}{\pi}\int_0^\pi \sin t\,\dd t
= \frac{1}{\pi}\bigl[-\cos t\bigr]_0^\pi = \frac{2}{\pi}
\approx 0.637 :
\]
فمتوسط قوس موجب كامل ليس $\frac12$ بل
$\frac2\pi$ --- إذ يقضي المنحنى وقتًا في الأعلى أكثر مما يقضيه
مثلث. وحسب \cref{exo:b1:integration:11} (مبرهنة التزايدات المنتهية
من أجل التكاملات، مع $g = 1$)، يكون المتوسط \emph{قيمة}: $\sin c =
\frac2\pi$ من أجل $c \in \intoo{0}{\pi}$ ما. وحسب \omterm{thm:b1:integration:riemann}{مجاميع ريمان}
في هذا الفصل، يكون المتوسط نهاية المتوسطات
العادية لعدد $n$ من العيّنات --- وهو الجسر بين المتوسط
المتقطّع للمعطيات والمتوسط المتصل لإشارة، وهو ما
يُدخل \omterm{thm:b1:integration:def}{التكامل} في الفيزياء.
\end{example}

\begin{theorem}[المكاملة بالتجزئة؛ التعويض]\label{thm:b1:integration:parts}
\begin{enumerate}
  \item إذا كانت $u, v$ من الصنف $C^1$ على $\intcc{a}{b}$:\index{مكاملة
        بالتجزئة}
        \[
        \int_a^b u'v = \bigl[uv\bigr]_a^b - \int_a^b uv' .
        \]
  \item إذا كانت $\varphi$ من الصنف $C^1$ على $\intcc{\alpha}{\beta}$ و $f$
        \omterm{def:b1:continuity:continuous}{متصلة} على $\varphi(\intcc{\alpha}{\beta})$:
        \index{تعويض}
        \[
        \int_{\alpha}^{\beta} f\bigl(\varphi(t)\bigr)\,\varphi'(t)\,
        \dd t = \int_{\varphi(\alpha)}^{\varphi(\beta)} f(x)\, \dd x .
        \]
\end{enumerate}
\end{theorem}

\begin{proof}
(1) $(uv)' = u'v + uv'$؛ فكامل على $\intcc{a}{b}$ وطبّق
المبرهنة الأساسية على الدالة $uv$ من الصنف $C^1$.

(2) لتكن $F$ دالة أصلية للمقدار $f$ على \omterm{prop:b1:reals:intervals}{فترة} الصورة
(\cref{thm:b1:integration:ftc}). عندئذ $(F \circ \varphi)' = (f
\circ \varphi)\,\varphi'$ (بقاعدة السلسلة)، ومنه يساوي الطرفان
$F(\varphi(\beta)) - F(\varphi(\alpha))$.
\end{proof}

\begin{example}\label{ex:b1:integration:parts}
$\int_0^1 t\,\eu^t \dd t = \bigl[t\,\eu^t\bigr]_0^1 - \int_0^1
\eu^t\dd t = \eu - (\eu - 1) = 1$. وبالتعويض $x =
\sin t$ ($t \in \intcc{0}{\frac\pi2}$):
\[
\int_0^1 \sqrt{1 - x^2}\, \dd x
= \int_0^{\pi/2} \cos t \cdot \cos t \, \dd t
= \int_0^{\pi/2} \frac{1 + \cos 2t}{2}\, \dd t = \frac\pi4 ,
\]
--- أي ربع القرص الواحدي، كما تقتضي الهندسة. (وهذا هو التخطيط
من \cref{met:b1:complex:trig} في العمل.)
\end{example}

\begin{remark}[مزالق شائعة في حساب التكامل]
(أ) \emph{يجب أن تكون التعويضات من الصنف $C^1$ على \omterm{prop:b1:reals:intervals}{الفترة} كلها}:
فالتغيير $x = \frac1t$ غير مشروع عبر $0$؛ وتطبيقه بلا تبصّر
على $\int_{-1}^{1}\frac{\dd x}{1 + x^2}$ «يبرهن» على أن
\omterm{thm:b1:integration:def}{التكامل} يساوي مقابله. وحين يكون للتعويض
شذوذ، فاقطع \omterm{prop:b1:reals:intervals}{الفترة} أولًا (شال)، ثم عوّض على
كل قطعة، ثم أعد التركيب. (ب) \emph{الدوال الأصلية
اللوغاريتمية تحتاج إلى قيم مطلقة}: $\int \frac{\dd x}{x - 2} =
\ln\abs{x - 2} + C$ على كل جانب من $2$ على حدة --- فكتابة
$\ln(x - 2)$ على $\intoo{0}{1}$ كتابةُ لوغاريتم عدد
سالب؛ وقد يختلف الثابت $C$ على جانبَي
الشذوذ. (ج) \emph{\omterm{thm:b1:integration:def}{التكامل} المنعدم لا
يقتل الدالة}: $\int_0^{2\pi}\sin = 0$؛ وتُشترط إيجابية
المكامَل قبل استنتاج $f = 0$
(\cref{ex:b1:integration:chasleswork}). (د) \emph{يجب معايرة \omterm{thm:b1:integration:riemann}{مجاميع
ريمان}}: ففي $\frac{b-a}{n}\sum f\bigl(a +
k\frac{b-a}{n}\bigr)$، على الخطوة خارجًا وعلى النقاط داخلًا
أن تطابقا التجزئة نفسها --- والخطأ الشائع مجموعٌ
$\sum_{k=1}^{n} f\bigl(\frac kn\bigr)$ \emph{دون} العامل
$\frac1n$، وهو يتباعد بدل أن يتقارب إلى $\int_0^1
f$. وقائمة التحقق قبل استدعاء \cref{thm:b1:integration:riemann}:
عمّل $\frac1n$ خارجًا، وأعد كتابة الحدّ في صورة $f$ عند $\frac kn$،
وسمِّ $f$ وتحقق من اتصالها.
\end{remark}

\begin{example}[خمّن، ثم اشتق، ثم عدّل]\label{ex:b1:integration:guessdiff}
ما هو $\int_1^x (\ln t)^2\,\dd t$؟ خمّن دالة أصلية من
الشكل $t\,P(\ln t)$ مع $P$ \omterm{def:b1:poly:def}{كثير حدود} واشتق:
\[
\bigl(t\,P(\ln t)\bigr)' = P(\ln t) + P'(\ln t) .
\]
ونحتاج إلى $P(u) + P'(u) = u^2$: فخذ $P(u) = u^2 - 2u + 2$
(بمطابقة المعاملات نزولًا من $u^2$). ومنه
\[
\int_1^x (\ln t)^2\,\dd t
= \bigl[t\bigl((\ln t)^2 - 2\ln t + 2\bigr)\bigr]_1^x
= x(\ln x)^2 - 2x\ln x + 2x - 2 ,
\]
وهي نتيجة تُبلَغ لولا ذلك بمكاملتين بالتجزئة. والفكرة
النافذة: من أجل المكامَلات من الشكل
(\omterm{def:b1:poly:def}{كثير حدود} في $\ln t$) أو (\omterm{def:b1:poly:def}{كثير حدود} مضروب في $\eu^{\lambda t}$)،
تكون للدالة الأصلية الشكل نفسه --- فاشتقاق تخمين مشكَّل
يحوّل المكاملة إلى جبر خطي على المعاملات،
أسرع وأقل عرضةً للخطأ من التجزئة المكرَّرة.
\end{example}

\begin{example}[التكامل البوميرانغ]\label{ex:b1:integration:boomerang}
احسب $I = \int_0^{\pi/2} \eu^x \cos x\,\dd x$. كامل
بالتجزئة مرتين، مشتقًا العامل المثلثي في كل مرة:
\[
I = \bigl[\eu^x \sin x\bigr]_0^{\pi/2}
- \int_0^{\pi/2} \eu^x \sin x\,\dd x
= \eu^{\pi/2} - J,
\qquad
J = \bigl[-\eu^x\cos x\bigr]_0^{\pi/2}
+ \int_0^{\pi/2} \eu^x\cos x\,\dd x = 1 + I .
\]
وقد عاد \omterm{thm:b1:integration:def}{التكامل} إلى نفسه: $I = \eu^{\pi/2} - 1 - I$،
ومنه
\[
I = \frac{\eu^{\pi/2} - 1}{2} .
\]
والفكرة النافذة: حين يكون المكامَل جداءَ دالتين
تعيدان إنتاج نفسيهما بالاشتقاق
($\eu^{ax}$ و $\cos bx$ و $\sin bx$)، تنتج مكاملتان بالتجزئة
معادلةً خطية في \omterm{thm:b1:integration:def}{التكامل} المجهول --- فحُلَّها
بدل أن تكامل؛ وبكيفية مكافئة، مرّ عبر
$\eu^{(a + \iu b)x}$ (\cref{ch:b1:complex}) وخذ الأجزاء
الحقيقية. والطريقان يعطيان الجواب نفسه، والتحقق من ذلك
اختبار سلامة مجاني.
\end{example}

\section{مجاميع ريمان}

\begin{theorem}[مجاميع ريمان]\label{thm:b1:integration:riemann}
لتكن $f$ \omterm{def:b1:continuity:continuous}{متصلة} على $\intcc{a}{b}$. عندئذ\index{مجاميع ريمان}
\[
S_n = \frac{b - a}{n} \sum_{k=0}^{n-1}
f\Bigl(a + k\,\frac{b-a}{n}\Bigr)
\xrightarrow[n \to \infty]{} \int_a^b f(t)\, \dd t ,
\]
وكذلك مع أيّ نقاط تقويم داخل الفترات الجزئية.
\end{theorem}

\begin{proof}
المقدار $S_n$ هو تكامل \omterm{def:b1:integration:step}{الدالة الدرجية} $\varphi_n$ المساوية
$f(a + k\frac{b-a}{n})$ على \omterm{prop:b1:reals:intervals}{الفترة} الجزئية $k$. ومن أجل
$\varepsilon > 0$ معطى، يعطي \omterm{def:b1:continuity:uniform}{الاتصال المنتظم} (هاينه) عددًا $\delta$؛ فمن أجل
$n > \frac{b-a}{\delta}$، تقع كل نقطة من \omterm{prop:b1:reals:intervals}{فترة} جزئية على بعد
$\delta$ من نقطة تقويمها، ومنه $\abs{f - \varphi_n} \leq
\varepsilon$ على $\intcc{a}{b}$، ومنه
\[
\Bigl| \int_a^b f - S_n \Bigr|
= \Bigl| \int_a^b (f - \varphi_n) \Bigr|
\leq (b-a)\,\varepsilon . \qedhere
\]
\end{proof}

\begin{omfigure}
\begin{tikzpicture}
\begin{axis}[omaxis, width=8.6cm, height=5.4cm,
    xmin=0, xmax=1.15, ymin=0, ymax=1.7,
    xtick={0,1}, ytick=\empty]
  \draw[fill=omDef!15, draw=omDef!60, thin]
    (axis cs:0,0)      rectangle (axis cs:0.125,0.4)
    (axis cs:0.125,0)  rectangle (axis cs:0.25,0.4156)
    (axis cs:0.25,0)   rectangle (axis cs:0.375,0.4625)
    (axis cs:0.375,0)  rectangle (axis cs:0.5,0.5406)
    (axis cs:0.5,0)    rectangle (axis cs:0.625,0.65)
    (axis cs:0.625,0)  rectangle (axis cs:0.75,0.7906)
    (axis cs:0.75,0)   rectangle (axis cs:0.875,0.9625)
    (axis cs:0.875,0)  rectangle (axis cs:1,1.1656);
  \addplot[omThm, very thick, domain=0:1.05, samples=60]
    {0.4 + x^2};
\end{axis}
\end{tikzpicture}

{\small مجموع ريمان يساري بعدد $n = 8$ مستطيلًا: فكلما تقلصت
الخطوة، فرض \omterm{def:b1:continuity:uniform}{الاتصال المنتظم} على مساحة السلّم أن تؤول نحو
$\int_a^b f$.}
\end{omfigure}

\begin{example}\label{ex:b1:integration:riemannexample}
$\displaystyle\sum_{k=1}^{n} \frac{1}{n + k} = \frac 1n
\sum_{k=1}^{n} \frac{1}{1 + k/n} \xrightarrow[n\to\infty]{}
\int_0^1 \frac{\dd x}{1 + x} = \ln 2$: أي نهاية غير مرئية
للحواصر الأولية، شفافة مجموعَ ريمان.
\end{example}

\begin{example}[مجموع ريمان ثانٍ، مع معايرة]\label{ex:b1:integration:riemann2}
جد $\lim_{n\to\infty} \sum_{k=1}^{n} \dfrac{n}{(n+k)^2}$.
عايِر:
\[
\sum_{k=1}^{n} \frac{n}{(n+k)^2}
= \frac{1}{n}\sum_{k=1}^{n} \frac{n^2}{(n+k)^2}
= \frac1n \sum_{k=1}^{n} \frac{1}{\bigl(1 + \frac kn\bigr)^2} ,
\]
وهو مجموع ريمان للدالة \omterm{def:b1:continuity:continuous}{المتصلة} $f(x) = \frac{1}{(1+x)^2}$ على
$\intcc{0}{1}$: فالنهاية هي
\[
\int_0^1 \frac{\dd x}{(1 + x)^2}
= \Bigl[-\frac{1}{1+x}\Bigr]_0^1 = \frac12 .
\]
والفكرة النافذة: كل الفنّ في السطر الأوسط ---
أي إجبار الحدّ على الشكل $f(\frac kn)$ بثمن
استخراج عامل $\frac1n$ واحد بالضبط؛ وما إن يستقم الشكل،
تقوم المبرهنة بالتحليل وتقوم المبرهنة الأساسية
بالحساب.
\end{example}

\begin{remark}[أين يعمل التكامل تاليًا]
لكل بناء من بناءات الفصل تتمّة. \omterm{thm:b1:integration:riemann}{فمجاميع ريمان}
تعود في \cref{ch:b1:series} جسرًا بين المتسلسلات و
التكاملات (بمقارنة $\sum \frac{1}{n^\alpha}$ بالمقدار $\int
\frac{\dd t}{t^\alpha}$)؛ والباقي التكاملي هو أدقّ
صيغ تايلور في \cref{ch:b1:taylor}؛ والتعريف
المبنيّ على $\sup$ هو النموذج الأوّلي لتكامل
لوبيغ في مجلّد السنة الثالثة، حيث تُعاد الخصائص الثلاث نفسها
(الخطية والرتابة ومبرهنة تقارب) على صنف
أوسع بكثير من الدوال. وتحوّل مسألة نهاية الأسبوع أدناه
المكاملةَ بالتجزئة إلى حساب أعداد: أي صمم
$\pi^2$.
\end{remark}

\section{تمارين}

\begin{exercise}[$\star$]\label{exo:b1:integration:1}
احسب: $\displaystyle\int_0^1 \frac{\dd x}{x^2 - 4}$
\emph{(بالعناصر البسيطة، \cref{ch:b1:fractions})}؛
$\displaystyle\int_1^{\eu} \ln t \,\dd t$؛
$\displaystyle\int_0^{\pi} t \sin t\, \dd t$.
\end{exercise}

\begin{exercise}[$\star$]\label{exo:b1:integration:2}
احسب $\displaystyle\int_0^{1} \frac{t}{(t^2+1)^2}\,\dd t$
(بالتعويض) و
$\displaystyle\int_0^{\pi/2} \cos^3 t\, \dd t$ (اكتب $\cos^3 =
\cos(1 - \sin^2)$).
\end{exercise}

\begin{exercise}[$\star$]\label{exo:b1:integration:3}
جد النهايات، مجاميعَ ريمان:
\[
u_n = \sum_{k=1}^{n} \frac{n}{n^2 + k^2},
\qquad
v_n = \frac{1}{n}\sqrt[n]{\frac{(2n)!}{n!\,n^n}}
\ \emph{(take logarithms)}.
\]
\end{exercise}

\begin{exercise}[$\star$]\label{exo:b1:integration:4}
لتكن $f$ \omterm{def:b1:continuity:continuous}{متصلة} على $\intcc{0}{1}$. احسب $\lim_{n\to\infty}
\int_0^1 x^n f(x)\,\dd x$. \emph{(اقطع $\intcc{0}{1}$ عند $1 -
\delta$.)}
\end{exercise}

\begin{exercise}[$\star\star$]\label{exo:b1:integration:5}
(كوشي--شوارتز) من أجل $f, g$ متصلتين على $\intcc{a}{b}$، برهن على
\[
\Bigl(\int_a^b fg\Bigr)^{\!2} \leq \int_a^b f^2 \cdot \int_a^b g^2 ,
\]
بنشر $\int_a^b (f + \lambda g)^2 \geq 0$ كثيرَ حدود من الدرجة الثانية في
$\lambda$. ومتى تكون مساواة؟
\end{exercise}

\begin{exercise}[$\star\star$]\label{exo:b1:integration:6}
لتكن $f$ \omterm{def:b1:continuity:continuous}{متصلة} على $\R$ ودورية بالدور $T$. برهن على أن $\int_a^{a+T}
f$ لا يتعلق بالمقدار $a$، وعلى أن $\frac1x \int_0^x f(t)\,\dd t
\to \frac 1T \int_0^T f$ عندما $x \to +\infty$.
\end{exercise}

\begin{exercise}[$\star\star$]\label{exo:b1:integration:7}
من أجل $f$ \omterm{def:b1:continuity:continuous}{متصلة} على $\intcc{0}{1}$ مع $\int_0^1 f = \frac12$،
برهن على أن للمقدار $f$ نقطةً صامدة في $\intcc{0}{1}$. \emph{(كامل
$f(x) - x$ واستعمل الإيجابية القطعية،
\cref{thm:b1:integration:props}~(4)، عبر نقيضها
مع مبرهنة القيم الوسطى.)}
\end{exercise}

\begin{exercise}[$\star\star\star$]\label{exo:b1:integration:8}
(تكاملات واليس) لتكن $W_n = \int_0^{\pi/2} \sin^n t\,\dd t$.
\begin{enumerate}
  \item برهن على العلاقة التراجعية $n W_n = (n-1) W_{n-2}$ ($n \geq 2$)
        بالتجزئة، واحسب $W_0, W_1$، ثم $W_{2p}$ و $W_{2p+1}$
        في صورة \omterm{def:b1:topology:closed}{مغلقة}.
  \item برهن على أن $(W_n)$ متناقصة مع $\frac{W_{n+1}}{W_n}
        \to 1$؛ وبرهن على أن المقدار $(n+1)\,W_{n+1} W_n$
        ثابت يساوي $\frac\pi2$؛ واستنتج المكافئ
        $W_n \sim \sqrt{\dfrac{\pi}{2n}}$.
\end{enumerate}
\end{exercise}

\begin{exercise}[$\star\star\star$]\label{exo:b1:integration:9}
(نيفن: العدد $\pi$ أصمّ) افترض $\pi = \frac ab$ مع $a, b \in
\N^*$، وضع، من أجل $n$ سيُختار،
\[
P(x) = \frac{x^n (a - bx)^n}{n!},
\qquad
I_n = \int_0^{\pi} P(x)\sin x\, \dd x .
\]
\begin{enumerate}
  \item برهن على أن $0 < I_n \leq \pi\,\frac{(\pi a)^n}{n!}$، وهو
        $< 1$ من أجل $n$ كبير.
  \item برهن على أن $P$ وكل مشتقاته تأخذ قيمًا \emph{صحيحة}
        عند $0$ وعند $\pi = \frac ab$. \emph{(بنشر ثنائي
        الحدّ: فمعاملات $P$ مضروبةً في $k!$ أعداد صحيحة
        من أجل $k \geq n$؛ و $P(\pi - x) = P(x)$.)}
  \item ضع $Q = P - P'' + P^{(4)} - \dots$ (وهو مجموع منته). تحقق
        من أن $\bigl(Q'\sin x - Q\cos x\bigr)' = P \sin x$، واستنتج
        أن $I_n = Q(\pi) + Q(0)$ عدد صحيح.
  \item اختم.
\end{enumerate}
\end{exercise}

\begin{exercise}[$\star\star\star$]\label{exo:b1:integration:10}
لتكن $f$ من الصنف $C^1$ على $\intcc{a}{b}$. برهن على النهاية من نوع
ريمان--لوبيغ
\[
\int_a^b f(t)\sin(\lambda t)\,\dd t \xrightarrow[\lambda \to
+\infty]{} 0
\]
\omterm{thm:b1:integration:parts}{بالمكاملة بالتجزئة}. ثم برهن عليها مرة أخرى من أجل $f$ \omterm{def:b1:continuity:continuous}{متصلة} فحسب،
بالتقريب المنتظم بالدوال الدرجية
(\cref{thm:b1:integration:approx}).
\end{exercise}

\begin{exercise}[$\star\star$]\label{exo:b1:integration:11}
(مبرهنة التزايدات المنتهية من أجل التكاملات) لتكن $f, g$ متصلتين على
$\intcc{a}{b}$ مع $g \geq 0$. برهن على أنه يوجد $c \in
\intcc{a}{b}$ يحقق
\[
\int_a^b f(t)\,g(t)\,\dd t = f(c)\int_a^b g(t)\,\dd t ,
\]
وبيّن بمثال أن الفرض $g \geq 0$ لا يمكن
إسقاطه.
\end{exercise}

\begin{exercise}[$\star\star\star$]\label{exo:b1:integration:12}
(العزوم تفرض الجذور) لتكن $f$ \omterm{def:b1:continuity:continuous}{متصلة} على $\intcc{a}{b}$
مع
\[
\int_a^b f(t)\,t^k\,\dd t = 0 \qquad \text{من أجل } k = 0, 1,
\dots, n .
\]
برهن على أن $f$ تنعدم عند $n + 1$ نقطة متمايزة من
$\intoo{a}{b}$. \emph{(إذا كانت $f$ لا تغيّر إشارتها إلا عند $z_1 < \dots
< z_m$ مع $m \leq n$، فكامل $f$ في مقابل $P(t) = (t - z_1)
\cdots (t - z_m)$ واستعمل الإيجابية القطعية.)}
\end{exercise}

\begin{remark}[منظورات داخل هذا المجلّد]
ثلاثة فصول قادمة تتّكئ على هذا الفصل مباشرة.
فالفصل \Cref{ch:b1:taylor} يحمل الباقي التكاملي --- إذ إن أدقّ
صيغ تايلور الثلاث مكاملةٌ بالتجزئة
مكرَّرة $n$ مرة. و \Cref{ch:b1:series} يحوّل تأطير
المجاميع بالتكاملات إلى المحك الحاسم من أجل $\sum n^{-\alpha}$،
وتلطّف مسألة نهاية الأسبوع فيه ذلك التأطير إلى ثابت أويلر.
و \Cref{ch:b1:curves} يجعل \omterm{thm:b1:integration:def}{التكامل} هندسيًا: إذ إن
طول قوس معلَّمي هو $\int \sqrt{x'(t)^2 +
y'(t)^2}\,\dd t$، وهو تكامل دالة \omterm{def:b1:continuity:continuous}{متصلة} على
قطعة --- أي الغرض المبنيّ هنا بالضبط، دون أيّ حاجة إلى
نظرية معتلّة. وأكثر الوقائع إعادةَ استعمال ستكون أشدَّها تواضعًا:
$\bigl|\int f\bigr| \leq (b - a)\sup\abs{f}$، وهي المتراجحة
التي تحوّل كل تقدير نقطي إلى تقدير تكاملي.
\end{remark}

\section{مسألة: آلة الصمم التكاملية}

\begin{problem}[{مسألة نهاية الأسبوع --- $\eu$ و $\pi^2$
أصمّان، و $\eu$ إلى ستة أرقام عشرية، و $\frac{22}{7} > \pi$ مع
برهان}]
\label{pb:b1:integration:1}
آلية واحدة تحرّك هذه المسألة كلها: \emph{لا يمكن أن يوجد مقدار
يجب أن يكون عددًا صحيحًا موجبًا ومع ذلك يُبرهن على أنه أصغر من $1$}.
وقد شغّل \Cref{exo:b1:integration:9} (نيفن) الآلية مرة واحدة
للبرهان على $\pi \notin \Q$؛ وهنا نصنّعها. وتحتاج الآلة
إلى ثلاثة أجزاء: مدخل \emph{صحيحية} (وهو قيم كثيرات
حدود محسنة الاختيار عند الطرفين)، ومدخل \emph{صغر} (وهو عامل
$\frac{1}{n!}$ يسحق \omterm{thm:b1:integration:def}{التكامل})، و\emph{جسر}
(هو \omterm{thm:b1:integration:parts}{المكاملة بالتجزئة}) يصل بينهما. ونبرهن على أن $\eu$
أصمّ ونحسبه بخطأ مصادَق عليه، ونبرهن على مبرهنة لوجاندر
الأدقّ التي تقول إن $\pi^2$ أصمّ، ونُختم بأبهج
تكامل في التحليل: $\int_0^1 \frac{x^4(1 -
x)^4}{1 + x^2}\dd x = \frac{22}{7} - \pi$، الذي يؤطّر $\pi$
باليد.\index{صمم!صمم $\pi^2$}\index{صمم!صمم $\eu$}

\medskip
\noindent\textbf{الجزء 1 --- الوقود.}
\begin{enumerate}
  \item برهن على أن $\dfrac{c^{\,n}}{n!} \to 0$ من أجل كل $c
        > 0$ مثبَّت \emph{(فبعد $n \geq 2c$، تنصّف كل خطوة الحدَّ
        على الأقل)}.
  \item (تكاملات بيتا) برهن، بالاستقراء على $m$ مع
        \omterm{thm:b1:integration:parts}{المكاملة بالتجزئة}:
        \[
        \int_0^1 x^{\,k}\,(1 - x)^{\,m}\,\dd x =
        \frac{k!\,m!}{(k + m + 1)!}
        \qquad (k, m \in \N).
        \]
  \item استنتج $\displaystyle\int_0^1 \bigl(x(1-x)\bigr)^n \dd x
        = \frac{1}{(2n+1)\binom{2n}{n}}$، و --- بالمقارنة مع
        الحاصر $x(1 - x) \leq \frac14$ --- التقدير
        $\binom{2n}{n} \geq \dfrac{4^n}{2n+1}$، المطابق
        للمقدار $\binom{2n}{n}^{1/n} \to 4$ من \cref{pb:b1:seq:1}.
  \item برهن على مبرهنة الصغر المساعدة المستعملة مرتين أدناه: من أجل كل
        دالة \omterm{def:b1:continuity:continuous}{متصلة} $g > 0$ على $\intoo{0}{1}$،
        \[
        0 < \int_0^1 \bigl(x(1-x)\bigr)^n g(x)\,\dd x
        \leq \frac{\sup_{\intcc{0}{1}}\abs g}{4^{\,n}} .
        \]
\end{enumerate}

\medskip
\noindent\textbf{الجزء 2 --- $\eu$: الصمم، ثم ستة
أرقام عشرية.} ضع $A_n = \displaystyle\int_0^1 x^n \eu^x \dd x$.
\begin{enumerate}[resume]
  \item احسب $A_0$ و $A_1$، وبرهن على العلاقة التراجعية $A_n = \eu
        - n A_{n-1}$، وعلى الحاصرين $0 < A_n \leq
        \dfrac{\eu}{n+1}$.
  \item بيّن بالاستقراء أن $A_n = \alpha_n + \beta_n \eu$
        مع $\alpha_n, \beta_n \in \Z$.
  \item استنتج أن $\eu$ أصمّ \emph{(فإذا كان $\eu = \frac
        pq$، فإن $q A_n$ عدد صحيح محبوس في
        $\intoo{0}{1}$ من أجل $n$ كبير)}. وقارن ببرهان
        \cref{exo:b1:seq:9}: النكتة نفسها، بوقود مختلف.
  \item برهن، بالاستقراء \omterm{thm:b1:integration:parts}{والمكاملة بالتجزئة}، على صيغة
        الباقي المضبوطة
        \[
        \eu = \sum_{k=0}^{n} \frac{1}{k!} + R_n,
        \qquad
        R_n = \frac{1}{n!}\int_0^1 (1 - t)^n\,\eu^{\,t}\,\dd t,
        \qquad
        \frac{1}{(n+1)!} \leq R_n \leq \frac{\eu}{(n+1)!} .
        \]
  \item خذ $n = 9$: حُدّ $R_9$ مستعملًا $\eu < 2.75$ (من
        $b_2 = 2.75$ في \cref{ex:b1:seq:e})، وقوّم المجموع،
        واختم بالتأطير المصادَق عليه $2.7182818 \leq
        \eu \leq 2.7182823$ --- أي ستة أرقام عشرية، $\eu \approx
        2.718282$، مع برهان.
\end{enumerate}

\medskip
\noindent\textbf{الجزء 3 --- مبرهنة لوجاندر: $\pi^2$
أصمّ.} لتكن $f(x) = \dfrac{x^n (1 - x)^n}{n!}$، ولنفترض
$\pi^2 = \frac ab$ مع $a, b \in \N^*$.
\begin{enumerate}[resume]
  \item بيّن $f(1 - x) = f(x)$ و $0 < f \leq
        \dfrac{1}{4^n\,n!}$ على $\intoo{0}{1}$.
  \item بيّن أن $f^{(k)}(0)$ و $f^{(k)}(1)$ عددان صحيحان من أجل
        كل $k \geq 0$ \emph{(انشر $x^n(1-x)^n$ بمعاملات
        صحيحة؛ و $\frac{k!}{n!} \in \Z$ من أجل $k \geq n$؛
        ثم استعمل التناظر)}.
  \item عرّف
        \[
        G = b^{\,n} \sum_{k=0}^{n} (-1)^k\,
        \pi^{2n - 2k} f^{(2k)} .
        \]
        بيّن أن $G(0)$ و $G(1)$ عددان صحيحان \emph{(لأن كل
        $b^n \pi^{2n-2k} = a^{\,n-k}\,b^{\,k}$)}.
  \item تحقق من التلسكوب $G'' + \pi^2 G = b^n \pi^{2n+2} f
        = \pi^2 a^n f$، ثم
        \[
        \frac{\dd}{\dd x}\Bigl(G'(x)\sin \pi x - \pi\,G(x)\cos
        \pi x\Bigr) = \pi^2 a^n f(x)\sin \pi x .
        \]
  \item كامل على $\intcc{0}{1}$ واستنتج
        \[
        \pi\,a^n \int_0^1 f(x)\sin(\pi x)\,\dd x = G(0) + G(1)
        \in \Z ,
        \]
        وهو عدد صحيح \emph{موجب} محدود بالمقدار $\dfrac{\pi
        a^n}{4^n\,n!}$.
  \item اختم مع السؤال 1 بأن $\pi^2$ أصمّ
        (لوجاندر، 1794)، وبأن هذا يقوّي
        \cref{exo:b1:integration:9}: فلماذا يستلزم صمم
        $\pi^2$ صممَ $\pi$، ولا يستلزم العكس؟
\end{enumerate}

\medskip
\noindent\textbf{الجزء 4 --- فهم الآلة.}
\begin{enumerate}[resume]
  \item حدّد القوتين المتعاكستين (صحيحية معطيات
        الطرفين؛ والصغر التحليلي للتكامل) و
        الجسرَ، في الجزأين 2 و 3. ثم فسّر لماذا يكون العامل
        $\frac{1}{n!}$ في $f$ هو العقدة: فإذا أُزيل،
        نجت الصحيحية، لكن أيّ متراجحة تموت، ومن أجل
        أيّ الكسور المدَّعاة $\frac ab$ يفشل البرهان عندئذ؟
  \item الفعّالية: افترض أن أحدهم يدّعي $\pi^2 = \frac ab$
        مع $a \leq 10$. بيّن أن التناقض يحطّ أصلًا
        عند $n = 7$: واحسب $\pi\,(10/4)^7/7! \approx
        0.38 < 1$. فالآلة لا تكتفي بالدحض؛ بل
        تدحض عند مرحلة مثبَّتة قابلة للحساب.
  \item تحقق من سلامة الجسر دون شرط: برهن بمكاملتين
        بالتجزئة على أن
        \[
        \int_0^1 x(1 - x)\sin(\pi x)\,\dd x = \frac{4}{\pi^3},
        \]
        ووفّق ذلك مع السؤال 14 عند $n = 1$ (وأبقِ
        $\pi^2$ رمزيًا: فالمتطابقة المتلسكبة تُقرأ $\pi^3
        \int_0^1 f_1 \sin \pi x = -(f_1''(0) + f_1''(1)) = 4$).
  \item ما الذي يجعل $\eu^x$ و $\sin \pi x$ مؤهلين نواتين
        للآلة؟ حدّد الخاصية (فكلٌّ منهما يحقق
        معادلة تفاضلية خطية بمعاملات
        ثابتة، ومنه تعود المكاملة المكرَّرة بالتجزئة
        إلى البداية)، وسمِّ الحدّ الأمامي: فالآلة نفسها،
        ملطَّفةً بيد إرميت ولينديمان، تبرهن على أن $\eu$
        و $\pi$ \emph{متساميان} --- وهذا خارج هذا المجلّد.
\end{enumerate}

\medskip
\noindent\textbf{الجزء 5 --- $\frac{22}{7}$ إزاء $\pi$، و
العبرة.}
\begin{enumerate}[resume]
  \item أثبت القسمة الكثيرحدودية
        \[
        \frac{x^4(1-x)^4}{1 + x^2}
        = x^6 - 4x^5 + 5x^4 - 4x^2 + 4 - \frac{4}{1 + x^2},
        \]
        واستنتج المتطابقة الشهيرة
        \[
        \int_0^1 \frac{x^4 (1-x)^4}{1 + x^2}\,\dd x
        = \frac{22}{7} - \pi .
        \]
  \item المكامَل موجب: فاستنتج $\pi <
        \frac{22}{7}$. ثم، بحصر $\frac{1}{1+x^2}$ بين
        $\frac12$ و $1$ وباستعمال $\int_0^1 (x(1-x))^4 =
        \frac{1}{630}$ (السؤال 3)، برهن على
        \[
        \frac{22}{7} - \frac{1}{630} \;\leq\; \pi \;\leq\;
        \frac{22}{7} - \frac{1}{1260} ,
        \]
        أي $3.14126 \leq \pi \leq 3.14207$: أي رقمان عشريان
        صحيحان، باليد.
  \item عمّم: بقسمة $x^{4m}(1-x)^{4m}$ على $1 + x^2$،
        بيّن أن الباقي هو الثابت $(-4)^m$ \emph{(اعمل
        بترديد $x^2 + 1$: $(1-x)^2 \equiv -2x$)}، واستنتج
        أعدادًا ناطقة $r_m$ تحقق
        \[
        \abs{\pi - r_m} \leq 4^{\,1 - 5m} ,
        \]
        وتحقق من أن $m = 1$ يستعيد السؤالين 20--21.
  \item قابل هذه الأعداد الناطقة بنظرية التقريب
        في \cref{pb:b1:derivative:1}: احسب $\abs{\pi -
        \frac{22}{7}} \approx 1.26\cdot10^{-3}$ إزاء
        ضمان ديريكليه $\frac{1}{49}$، واذكر
        $\abs{\pi - \frac{355}{113}} \approx 2.7\cdot10^{-7}$
        إزاء $\frac{1}{113^2} \approx 7.8\cdot10^{-5}$:
        فتوجد من أجل $\pi$ تقريبات ناطقة جيدة استثنائيًا
        --- وهذا متسق، إذ ليس معروفًا أن $\pi$ سيّئ
        التقريب.
  \item (فخّ العدد الصحيح، مجرَّدًا) برهن على المبرهنة المساعدة التي
        توحّد كل شيء: \emph{إذا كان $x \in \R$ ووُجدت
        أعداد صحيحة $a_n, b_n$ تحقق $0 < \abs{a_n + b_n x} \to 0$،
        فإن $x$ أصمّ.} واذكر أمثلتها في هذه
        المسألة، وفي \cref{exo:b1:integration:9}، وفي
        \cref{exo:b1:seq:9}، وفي \cref{pb:b1:derivative:1}.
  \item توليفة، جملة واحدة لكلٍّ: (أ) أجزاء الآلة الثلاثة
        وأين يعيش كلٌّ منها في عدّة هذا
        الفصل؛ (ب) وما الذي يسهم به \omterm{thm:b1:integration:def}{التكامل} ولم تستطع
        مبرهنة التزايدات المنتهية في \cref{pb:b1:derivative:1}
        أن تسهم به؛ (ج) وجرد النتائج المستخرجة (صممان
        وثابت من ستة أرقام عشرية وتأطير
        واحد للمقدار $\pi$ وحاصر ثنائي حدّ واحد)؛ (د) و
        الحدّ الأمامي (إرميت ولينديمان؛ والفخّ نفسه، مشغَّلًا على
        $\zeta(2)$ و $\zeta(3)$، في حساب الأعداد في القرن
        العشرين).
\end{enumerate}
\end{problem}
